\chapter{État de l'art et positionnement}
\label{chapitre:etatart}

Le travail se situe dans la littérature existante. Les approches automatisées de
détection sont passées en revue, puis les deux méthodes non supervisées retenues comme
comparateurs. Viennent ensuite la détection temporelle de changements persistants, dont le
système proposé s'inspire, et l'hétérogénéité des protocoles qui limite la comparaison des
performances publiées. Le positionnement du mémoire clôt le chapitre.

% =====================================================================
\section{Approches automatisées de détection}
% =====================================================================

Les approches fondées sur des capteurs cherchent à repérer des variations de l'activité, des
transitions posturales ou de la dynamique du mouvement susceptibles d'être associées à la
boiterie, sans être spécifiques à cette affection. La revue de la littérature fait ressortir
trois grandes familles: l'analyse de la démarche par vision par ordinateur
\citep{vanhertem2014evaluation,myint2024lameness}, la classification de comportements ou
d'états locomoteurs à partir de données accélérométriques
\citep{oleary2020accelerometers,alsaaod2019automatic,balasso2023behaviour,
benaissa2019calving,lavrova2023lameness}, et la détection d'anomalies sur des indicateurs
dérivés \citep{qiao2021review,chandola2009anomaly}. Les synthèses disponibles soulignent la
prédominance des approches supervisées et le manque d'études non supervisées conduites dans
des conditions réelles d'élevage \citep{qiao2021review,oleary2020accelerometers}. Une revue
consacrée à l'intelligence artificielle appliquée à la détection automatisée relève également
trois freins persistants: l'hétérogénéité des protocoles, la dépendance aux étiquettes et le
faible nombre d'évaluations en conditions réelles \citep{siachos2024automated}.

À titre d'exemple, \citet{vanhertem2014evaluation} exploitent des enregistrements vidéo 3D
successifs, tandis que \citet{myint2024lameness} proposent un système en temps réel fondé sur
une caméra latérale. Ces méthodes obtiennent de bonnes performances, mais nécessitent une
infrastructure de vision et des données étiquetées, deux conditions qui ne sont pas réunies
dans le corpus IceQube non étiqueté. Des jeux de données publics annotés pour la boiterie
commencent aussi à apparaître, comme \textit{CowScreeningDB}
\citep{ismail2024cowscreeningdb}; ils restent toutefois rares, liés à leur propre dispositif
d'acquisition et distincts du corpus exploité dans ce mémoire.

Ces approches peuvent également être comparées selon trois dimensions complémentaires: la
source des données, la stratégie analytique et la nature de la sortie
\citep{qiao2021review,oleary2020accelerometers}. Ce cadre évite de juxtaposer des systèmes qui
ne répondent pas au même besoin. Un modèle supervisé fondé sur des images apprend à attribuer
une classe à partir de données étiquetées, alors qu'un système d'alerte fondé sur des capteurs
d'activité signale une déviation comportementale qui doit être vérifiée. Les deux systèmes
diffèrent donc par les données requises, le degré d'interprétabilité de leur sortie et leur
usage en ferme.

La Figure~\ref{fig:taxonomie_approches} synthétise cette organisation des approches
automatisées. Le premier axe porte sur la source des données, depuis l'observation visuelle
jusqu'aux accéléromètres portés. Le deuxième décrit la stratégie analytique, qui peut reposer
sur des règles et des seuils, un apprentissage supervisé, une détection non supervisée
d'anomalies ou une approche hybride. Le troisième précise la nature de la sortie, du score
locomoteur à l'alerte nécessitant une vérification. Le présent travail relève de la
\textbf{branche non supervisée et hybride}: il combine une détection de changements
comportementaux sans étiquettes avec des règles temporelles et opérationnelles. Cette branche
est encadrée en vert sur la figure.

\begin{figure}[H]
    \centering
    \resizebox{\textwidth}{!}{%
    \begin{tikzpicture}[
        block/.style={rectangle, draw=black!70, rounded corners=3pt, minimum width=4cm,
                      minimum height=1cm, align=center, font=\small\bfseries, inner sep=8pt},
        small/.style={rectangle, draw=black!35, rounded corners=2pt, minimum width=2.8cm,
                      minimum height=0.85cm, text width=2.6cm, align=center, font=\footnotesize, inner sep=6pt},
        arrow/.style={-{Stealth[length=2.5mm]}, thick, black!55}
    ]
        \node[block, fill=ifblue!12] (data) at (0,0) {Sources de données};
        \node[block, fill=loforange!12] (model) at (7,0) {Stratégie analytique};
        \node[block, fill=rulegreen!12] (output) at (14,0) {Sortie opérationnelle};

        \node[small, fill=ifblue!6] (obs) at (-1.8,-2.2) {Observation\\et notation};
        \node[small, fill=ifblue!6] (vision) at (1.8,-2.2) {Vision\\de la démarche};
        \node[small, fill=ifblue!6] (accel) at (-1.8,-4.0) {Accéléromètres\\activité / couchage};
        \node[small, fill=ifblue!6] (multi) at (1.8,-4.0) {Multi-source\\capteurs + contexte};

        \node[small, fill=loforange!6] (rules) at (5.2,-2.2) {Règles\\et seuils};
        \node[small, fill=loforange!6] (sup) at (8.8,-2.2) {Supervisé};
        \node[small, fill=loforange!6] (unsup) at (5.2,-4.0) {Non supervisé\\(anomalies)};
        \node[small, fill=loforange!6] (hyb) at (8.8,-4.0) {Hybride\\score + règles};

        \node[small, fill=rulegreen!6] (score) at (12.2,-2.2) {Score\\locomoteur};
        \node[small, fill=rulegreen!6] (classif) at (15.8,-2.2) {Classe binaire\\boiteux / non};
        \node[small, fill=rulegreen!6] (episode) at (12.2,-4.0) {Épisode\\d'alerte};
        \node[small, fill=rulegreen!6] (alerte) at (15.8,-4.0) {Alerte de\\vérification};

        \draw[arrow] (data) -- (model);
        \draw[arrow] (model) -- (output);

        % Lignes: parent -> rang 1 (directes, courtes)
        \draw[black!30] (data.south) -- ++(0,-0.4) -| (obs.north);
        \draw[black!30] (data.south) -- ++(0,-0.4) -| (vision.north);
        \draw[black!30] (model.south) -- ++(0,-0.4) -| (rules.north);
        \draw[black!30] (model.south) -- ++(0,-0.4) -| (sup.north);
        \draw[black!30] (output.south) -- ++(0,-0.4) -| (score.north);
        \draw[black!30] (output.south) -- ++(0,-0.4) -| (classif.north);

        % Lignes: rang 1 -> rang 2 (verticales simples)
        \draw[black!30] (obs.south) -- (accel.north);
        \draw[black!30] (vision.south) -- (multi.north);
        \draw[black!30] (rules.south) -- (unsup.north);
        \draw[black!30] (sup.south) -- (hyb.north);
        \draw[black!30] (score.south) -- (episode.north);
        \draw[black!30] (classif.south) -- (alerte.north);

        % Encadrement du positionnement de ce mémoire (autour des nœuds unsup et hyb)
        \node[draw=rulegreen, thick, rounded corners=3pt, inner sep=6pt,
              fit=(unsup)(hyb)] (encadre) {};
        \node[font=\footnotesize\bfseries, rulegreen!80!black, below=4pt of encadre]
            {Positionnement de ce mémoire};
    \end{tikzpicture}%
    }
    \caption{Taxonomie des approches de détection automatisée.}
    \label{fig:taxonomie_approches}
\end{figure}

Du point de vue de la stratégie analytique, les méthodes de détection se déclinent
classiquement en deux familles:
\begin{itemize}
    \item[(a)] \textit{Modèles supervisés}: lorsque des étiquettes fiables et suffisamment
    représentatives sont disponibles. Ces approches peuvent atteindre de bonnes performances
    de classification
    \citep{vanhertem2014evaluation,viazzi2013analysis}, mais leur
    dépendance aux étiquettes limite leur déployabilité dans des contextes où l'étiquetage est
    incomplet ou incertain;
    \item[(b)] \textit{Détection d'anomalies et règles expertes}: lorsque les étiquettes sont
    partielles ou absentes. Ces approches peuvent être utilisées sans exemples positifs pour
    l'entraînement, notamment lorsque le suivi locomoteur n'est pas systématiquement
    synchronisé avec les données capteurs
    \citep{qiao2021review,oleary2020accelerometers}.
\end{itemize}

Le corpus IceQube décrit à la Section~\ref{sec:donnees} ne contient pas d'étiquettes cliniques
synchronisées permettant un apprentissage supervisé. Cette contrainte motive le choix d'une
architecture qui combine des mesures de déviation et des règles temporelles explicites
\citep{chandola2009anomaly,palma2025scoping}. Dans un autre contexte bovin,
\citet{michelena2025comparative} compare plusieurs détecteurs non supervisés pour repérer les
chaleurs. Cette étude ne fournit pas de preuve sur la boiterie, mais elle montre l'intérêt
d'une comparaison empirique entre plusieurs détecteurs appliqués aux mêmes données. Cette
distinction conduit à retenir Isolation Forest et LOF comme références non supervisées afin
d'évaluer l'apport de la structure temporelle explicite du système proposé.

% =====================================================================
\section{Méthodes non supervisées de référence}
% =====================================================================

\subsection{Isolation Forest}
La première méthode de référence repose sur l'algorithme Isolation Forest (IF), proposé par
\citet{liu2008isolation}. Son principe est que les anomalies sont plus faciles à isoler que
les points normaux dans un
espace de variables. Son échantillonnage et ses arbres aléatoires permettent de traiter des
jeux tabulaires sans modèle explicite de la distribution normale. Son implémentation dans
scikit-learn, comme montré dans \citet{sklearn_api}, fournit un comparateur ponctuel
reproductible; elle ne modélise toutefois pas directement la direction ni la durée d'un
changement.

La Figure~\ref{fig:if_principle} illustre le principe fondamental d'IF. Une anomalie
(point rouge) est isolée en seulement deux coupures aléatoires, tandis que les points
normaux (cluster bleu) nécessitent de nombreuses partitions supplémentaires. Le score
d'anomalie est inversement proportionnel au nombre de coupures.

\begin{figure}[H]
    \centering
    \resizebox{0.72\textwidth}{!}{%
    \begin{tikzpicture}
        % Cadre
        \draw[black!50] (0,0) rectangle (12,6);

        % --- Cluster normal : 15 points groupés à gauche ---
        \foreach \x/\y in {2.0/3.0, 2.3/3.4, 2.6/2.8, 2.2/3.2, 1.9/3.5,
                           2.5/3.0, 2.1/2.7, 2.4/3.3, 2.7/3.5, 1.8/3.1,
                           2.2/2.9, 2.5/3.4, 2.3/2.8, 2.1/3.3, 2.6/3.2} {
            \fill[ifblue] (\x,\y) circle (2.5pt);
        }

        % --- Anomalie : 1 point isolé à droite ---
        \fill[rawred] (9.5,4.5) circle (4pt);

        % --- Coupure 1 : verticale entre cluster et anomalie ---
        \draw[black!70, dashed, thick] (6,0.3) -- (6,5.7);

        % --- Coupure 2 : horizontale, zone droite seulement ---
        \draw[black!70, dashed, thick] (6.3,2.8) -- (11.7,2.8);

        % --- Étiquettes hors figure (marges) ---
        % Cluster
        \node[font=\footnotesize, black] at (2.3,1.5) {Points normaux};
        \node[font=\footnotesize, black!70] at (2.3,1.0) {difficiles à isoler};

        % Anomalie
        \node[font=\footnotesize, black] at (9.5,5.3) {Anomalie};
        \node[font=\footnotesize, black!70] at (9.5,3.7) {isolée en 2 coupures};

        % Numéros des coupures
        \node[font=\footnotesize\bfseries, black, fill=white, inner sep=2pt] at (6,0.9) {1};
        \node[font=\footnotesize\bfseries, black, fill=white, inner sep=2pt] at (10.8,2.8) {2};
    \end{tikzpicture}%
    }
    \caption{Principe d'Isolation Forest.}
    \label{fig:if_principle}
\end{figure}

\subsection{Local Outlier Factor}
La méthode nommée Local Outlier Factor (LOF), proposée par
\citet{breunig2000lof}, estime la
densité locale de chaque point par rapport à celle de ses voisins. Il vise ainsi des points
rares dans leur voisinage, tandis qu'IF repose sur leur facilité d'isolation globale. Le coût
de LOF dépend notamment de la recherche des voisins et de la dimension des données. Dans ce
mémoire, il sert de second comparateur non supervisé; aucune supériorité générale ne lui est
attribuée a priori.

La Figure~\ref{fig:lof_principle} illustre ce principe. Un point situé au cœur du groupe a des
voisins proches, donc une densité locale comparable à la leur. Le point isolé conserve un
voisinage, mais nettement plus étendu~: sa densité locale est faible devant celle de ses propres
voisins, et c'est ce rapport, et non une distance absolue, qui détermine son score.

\begin{figure}[H]
    \centering
    \resizebox{0.72\textwidth}{!}{%
    \begin{tikzpicture}
        % Cadre
        \draw[black!50] (0,0) rectangle (12,6);

        % --- Groupe dense a gauche ---
        \foreach \x/\y in {2.0/3.0, 2.3/3.4, 2.6/2.8, 2.2/3.2, 1.9/3.5,
                           2.5/3.0, 2.1/2.7, 2.4/3.3, 2.7/3.5, 1.8/3.1,
                           2.2/2.9, 2.5/3.4, 2.3/2.8, 2.1/3.3, 2.6/3.2} {
            \fill[ifblue] (\x,\y) circle (2.5pt);
        }

        % --- Voisinage d'un point du groupe : rayon court ---
        \draw[black!70, dashed] (2.35,3.1) circle (0.85);

        % --- Points clairsemes a droite ---
        \foreach \x/\y in {7.9/3.0, 10.4/4.6, 8.3/4.9, 10.6/2.6} {
            \fill[ifblue] (\x,\y) circle (2.5pt);
        }

        % --- Point isole : voisinage bien plus etendu ---
        \fill[rawred] (9.2,3.8) circle (4pt);
        \draw[black!70, dashed] (9.2,3.8) circle (1.6);

        % --- Etiquettes ---
        \node[font=\footnotesize, black] at (2.35,1.6) {Voisinage dense};
        \node[font=\footnotesize, black!70] at (2.35,1.1) {densité proche de celle des voisins};

        \node[font=\footnotesize, black] at (9.2,5.7) {Voisinage clairsemé};
        \node[font=\footnotesize, black!70] at (9.2,1.7) {densité faible devant celle des voisins};
    \end{tikzpicture}%
    }
    \caption{Principe de Local Outlier Factor.}
    \label{fig:lof_principle}
\end{figure}

\subsection{Comparaison synthétique IF et LOF}
Le Tableau~\ref{tab:if_vs_lof} synthétise les différences fondamentales entre les deux
détecteurs utilisés dans ce projet.

\begin{table}[H]
    \centering
    \caption{Comparaison des propriétés d'Isolation Forest et de Local Outlier Factor.}
    \label{tab:if_vs_lof}
    \small
    \begin{tabular}{>{\raggedright\arraybackslash}p{3.6cm}
                    >{\raggedright\arraybackslash}p{5.8cm}
                    >{\raggedright\arraybackslash}p{5.8cm}}
        \toprule
        \textbf{Propriété}            & \textbf{Isolation Forest}
            & \textbf{Local Outlier Factor} \\
        \midrule
        Principe                      & Isolabilité par partitions
            & Densité locale relative \\
        Type d'anomalie ciblée        & Globale (isolée dans l'espace)
            & Locale (faible densité relative) \\
        Référence fondatrice          & \citet{liu2008isolation}
            & \citet{breunig2000lof} \\
        Dépendance principale         & Nombre d'arbres et contamination
            & Nombre de voisins et contamination \\
        Limite pour ce projet         & Ne modélise pas la persistance
            & Ne modélise pas la persistance \\
        Paramètres critiques          & \texttt{contamination}, \texttt{n\_estimators}
            & \texttt{n\_neighbors}, \texttt{contamination} \\
        \bottomrule
    \end{tabular}
\end{table}

IF fournit un comparateur fondé sur l'isolation globale, tandis que LOF apporte un contraste
fondé sur la densité locale. La distinction apparaît sur les cas où les deux divergent. IF
signale un intervalle globalement atypique, même lorsqu'il occupe une région peu peuplée de
l'espace des variables, alors que LOF retient un intervalle ordinaire dans l'absolu mais
s'écartant de son voisinage immédiat. Aucun des deux ne tient toutefois compte de l'ordre des
intervalles~: chacun juge un instant sans le situer dans la trajectoire de la vache. La sortie
principale du mémoire repose donc sur les détecteurs temporels décrits au
Chapitre~\ref{chapitre:systeme}, et la section suivante examine les méthodes qui exploitent
cette dimension temporelle.

% =====================================================================
\section{Détection temporelle de changements persistants}
% =====================================================================

La boiterie n'est pas nécessairement représentée par un point extrême isolé. Les changements
rapportés dans les signaux obtenus concernent souvent une tendance, une rupture ou une
modification persistante par rapport au comportement habituel de l'animal
\citep{alsaaod2019automatic,paudyal2021lying}. Cette structure temporelle favorise les
méthodes capables d'accumuler une évidence unilatérale plutôt que de classer indépendamment
chaque intervalle.
La méthode des sommes cumulées \citep{page1954continuous} constitue une famille
classique de détection de changement, formalisée depuis en analyse séquentielle et en
détection de rupture \citep{basseville1993detection} et toujours au cœur des méthodes
actuelles \citep{truong2020selective}. Appliquées à une combinaison de déficits normalisés,
elles permettent d'oublier les fluctuations brèves, d'accumuler une baisse soutenue et de
déclencher une alerte lorsque l'évidence franchit un seuil. Le problème traité n'est donc pas
seulement une classification ponctuelle, mais la décision progressive qu'une trajectoire
individuelle s'éloigne durablement de sa référence. Dans un contexte de production bovine,
cette logique doit
être associée à une période de référence individuelle, au cycle journalier, à la qualité des
données et à la concordance de plusieurs familles de signaux. Elle reste toutefois une détection
comportementale: seule une validation avec examens locomoteurs et diagnostics contemporains
peut établir la relation clinique. Encore faut-il que les travaux publiés soient comparables,
ce que la section suivante met en doute.

% =====================================================================
\section{Hétérogénéité des protocoles~: un obstacle à la comparaison}
\label{sec:heterogeneite}
% =====================================================================

Les revues spécialisées sur la détection automatisée de boiterie convergent vers une
limite récurrente: l'hétérogénéité des protocoles. Le type de capteur, sa position sur
l'animal, la définition opérationnelle de la boiterie, la taille des cohortes et l'horizon
d'évaluation varient fortement d'une étude à l'autre, ce qui complique la comparaison
directe des performances publiées, comme le soulignent les travaux présentés dans
\citet{oleary2020accelerometers} et \citet{alsaaod2019automatic}.

Les travaux représentatifs du Tableau~\ref{tab:revue_ciblee} montrent pourquoi les
performances numériques ne doivent pas être juxtaposées sans précaution: les définitions
cliniques, les populations, les modalités de capteurs, les découpages d'évaluation et les
unités statistiques diffèrent. La présente section examine donc les protocoles plutôt qu'un
classement de sensibilités.

\begin{table}[H]
    \centering
    \small
    \caption{Travaux directement reliés aux choix méthodologiques.}
    \label{tab:revue_ciblee}
    \begin{tabular}{>{\raggedright\arraybackslash}p{3.0cm}
                    >{\raggedright\arraybackslash}p{4.7cm}
                    >{\raggedright\arraybackslash}p{6.0cm}}
        \toprule
        \textbf{Travail} & \textbf{Constat utile} & \textbf{Conséquence pour ce mémoire} \\
        \midrule
        \citet{alsaaod2019automatic}
        & Les systèmes automatiques reposent sur des modalités et des protocoles hétérogènes.
        & Limiter la sortie à une alerte comportementale et documenter le protocole complet. \\
        \addlinespace
        \citet{oleary2020accelerometers}
        & Les variables accélérométriques associées à la boiterie varient selon le capteur et
          le protocole.
        & Éviter une signature universelle et conserver plusieurs familles de signaux. \\
        \addlinespace
        \citet{lavrova2023lameness}
        & Une étude multi-fermes combine activité, couchage, production et caractéristiques
          individuelles.
        & Comparer chaque vache à sa propre référence et reconnaître l'intérêt futur de
          variables contextuelles. \\
        \addlinespace
        \citet{balasso2023behaviour}
        & Les signaux accélérométriques bruts permettent de modéliser des motifs
          comportementaux complexes.
        & Ne pas transposer ces capacités à des agrégats de 15~minutes qui ne contiennent
          plus la dynamique fine de la démarche. \\
        \addlinespace
        \citet{palma2025scoping}
        & Le déploiement de l'intelligence artificielle en élevage dépend de la qualité des
          données, de l'intégration et de l'usage.
        & Privilégier une sortie traçable, interprétable et destinée à une vérification
          humaine. \\
        \bottomrule
    \end{tabular}
\end{table}
Le travail présenté dans \citet{oleary2020accelerometers} sur les accéléromètres pour la
détection de boiterie confirme cette dispersion et souligne l'absence de standardisation
dans les protocoles d'évaluation.
\citet{alsaaod2019automatic} soulignent aussi en particulier l'absence
de véritables études de longue durée ayant abouti à un système d'alerte précoce valide de
bout en bout.
Une étude longitudinale menée dans six fermes allemandes combine activité, couchage,
production laitière et caractéristiques individuelles dans un modèle à effets mixtes
\citep{lavrova2023lameness}. Elle illustre l'intérêt d'informations hétérogènes, mais son
protocole supervisé et ses données contextuelles diffèrent de ceux du présent mémoire.

Enfin, la qualité de la mesure amont reste déterminante: l'échantillonnage, l'édition des
épisodes et la représentation des comportements influencent directement la fiabilité des
variables dérivées \citep{ledgerwood2010lying,balasso2023behaviour}.

Concrètement, l'hétérogénéité des protocoles se manifeste sur au moins cinq dimensions:
\begin{enumerate}
    \item \textit{La définition du cas}: classe clinique, score locomoteur seuil ou probabilité
    de boiterie selon les études, d'où des phénomènes mesurés et des sensibilités
    difficilement comparables;
    \item \textit{L'instrumentation}: position du capteur (patte
    \citep{alsaaod2012electronic,thorup2015lameness}, collier
    \citep{vanhertem2013lameness}, vision dorsale \citep{viazzi2013analysis}), fréquence
    d'échantillonnage et signaux
    contextuels varient, produisant des variables de nature différente;
    \item \textit{La granularité analytique}: prédiction instantanée, journalière, par épisode
    ou par vache, avec des exigences de précision temporelle très différentes;
    \item \textit{L'objectif métier}: classification clinique, alerte précoce, filtrage ou
    score de risque, chaque objectif imposant des coûts d'erreur différents;
    \item \textit{L'évaluation}: sensibilité, spécificité, AUC, taux d'alertes, IoU ou
    robustesse inter-fermes, un même système paraissant excellent ou médiocre selon la
    métrique.
\end{enumerate}

Autrement dit, deux articles peuvent rapporter une « bonne performance » tout en répondant à
des problèmes différents. Cette observation justifie le choix de consacrer un chapitre entier
aux protocoles d'évaluation et à leurs artefacts, plutôt que de se limiter à un score unique
agrégé.

\FloatBarrier

% =====================================================================
\section{Synthèse critique et positionnement du mémoire}
% =====================================================================

\subsection{Choix méthodologique et travaux directement pertinents}
\label{sec:choix_methodo}
La revue de littérature présentée est narrative et ciblée; elle ne constitue ni une revue systématique
ni une méta-analyse. Les travaux ont été retenus parce qu'ils éclairent directement l'une
des décisions du projet: nature indirecte des signaux comportementaux, absence de labels
cliniques dans le corpus principal, intérêt d'une référence individuelle, nécessité d'une
évaluation temporelle et exigence d'interprétabilité.

Les modèles supervisés et les réseaux profonds peuvent exploiter des motifs complexes lorsque
des signaux suffisamment riches et des labels contemporains sont disponibles
\citep{vanhertem2014evaluation,balasso2023behaviour}. Ces conditions ne sont pas réunies dans ce mémoire: les variables IceQube sont déjà agrégées par intervalles de 15~minutes et le corpus
principal ne comporte pas de diagnostic locomoteur synchronisé. Une classification clinique
supervisée ne serait donc pas défendable. Les méthodes non supervisées restent utiles comme
points de comparaison, mais elles ne résolvent pas à elles seules la direction ni la
persistance d'un changement. Des travaux récents transposent la détection d'anomalies non
supervisée à des séries de capteurs d'élevage bruitées \citep{kakar2024timetector}, ce qui
conforte l'intérêt d'approches sans étiquettes; ils restent toutefois orientés vers la
reconstruction de signaux plutôt que vers une baisse comportementale directionnelle,
persistante et individualisée comme celle que cible HYPO.

Le Tableau~\ref{tab:revue_ciblee}, présenté à la section~\ref{sec:heterogeneite}, réunit
les travaux les plus directement reliés à ces choix méthodologiques. Il ne compare pas leurs
performances numériques, car les capteurs, les définitions de cas et les unités statistiques
diffèrent.

Ces travaux motivent une architecture classique, individualisée et temporelle. IF et LOF
sont conservés comme comparateurs ponctuels de familles différentes; le résultat principal
repose toutefois sur une accumulation directionnelle des baisses et sur la concordance de
plusieurs signaux. L'apprentissage profond n'est pas rejeté en général: il est simplement
hors de portée des données agrégées et non étiquetées utilisées dans ce mémoire.

\subsection{Synthèse critique}
La littérature et les guides pratiques indiquent trois exigences pour une solution de
surveillance associée au risque de boiterie:
\begin{enumerate}
  \item \textit{Réactivité opérationnelle}: repérer des changements persistants assez tôt
  pour prioriser l'observation;
  \item \textit{Maîtrise de la charge d'alerte}: limiter les notifications sans suite opérationnelle
  et la fatigue d'alerte, qui dégradent rapidement la confiance des utilisateurs
  \citep{rutten2013invited};
  \item \textit{Interprétabilité}: fournir des raisons explicites utiles à la vérification pour
  chaque alerte, afin de faciliter la prise de décision terrain.
\end{enumerate}

En pratique, les méthodes purement pilotées par les données peuvent manquer d'explicabilité pour
l'opérationnel, tandis que les approches purement à règles peuvent manquer de flexibilité
face à la variabilité inter-animaux. Un compromis robuste consiste à coupler la détection
d'anomalies et des contraintes métier explicites. Ce couplage cumule la
capacité de généralisation des algorithmes non supervisés et des gardes-fous
interprétables. Des travaux récents confirment cette orientation: les systèmes
d'élevage de précision les plus prometteurs combinent des signaux hétérogènes avec des
logiques de décision interprétables, plutôt que de s'appuyer sur un modèle unique en boîte
noire \citep{palma2025scoping,michelena2025plf_review}.

L'hétérogénéité des protocoles examinée à la Section~\ref{sec:heterogeneite} a une conséquence
directe sur l'évaluation des systèmes: de nombreux travaux comparent des capteurs, des fermes, des
labels et des définitions de cas différents, ce qui complique l'appréciation d'une chaîne
déployable de bout en bout. Dans ce contexte, il existe une place claire pour des
contributions moins théoriques mais plus systémiques: \textbf{une chaîne reproductible, une
évaluation explicite sur plusieurs niveaux et une traçabilité des artefacts expérimentaux}.

La revue de littérature conduit alors à un enseignement central: la difficulté n'est pas
seulement de proposer un détecteur performant, mais de construire une chaîne de décision
cohérente, reproductible et défendable du point de vue «~métier~». Une telle exigence justifie
la structure du présent mémoire, qui traite l'algorithme, les règles métier, l'interface et
l'évaluation comme un ensemble cohérent plutôt que comme des blocs séparés.

\subsection{Positionnement du mémoire}
Le présent mémoire retient d'abord une approche temporelle directionnelle orientée vers les
baisses d'activité, puis examine une extension bidirectionnelle. Cette hiérarchie reflète le
niveau de preuve: la branche \textbf{HYPO} est soumise à une évaluation attribuable, une ablation et
une analyse de sensibilité complètes, puis la branche \textbf{d'instabilité} teste une
hypothèse complémentaire qui reste exploratoire. La contribution relève de l'ingénierie scientifique:
construire une chaîne reproductible qui compare chaque vache à sa propre période de référence et évaluer
honnêtement ce que permettent des données d'activité non étiquetées. Concrètement, le
mémoire livre:
\begin{enumerate}
  \item une \textbf{branche HYPO} évaluée sur 44 événements postérieurs à la période de référence, avec exécution propre de
  référence, ablation contre IF et LOF, analyse OFAT de 20 variantes et stress test
  à aire d'enveloppe appariée sur 198 événements;
  \item une \textbf{branche d'instabilité} exploratoire et cinq stratégies de fusion, dont une règle
  hiérarchique séparant surveillance et notification de vérification;
  \item une \textbf{campagne étendue} de 132 événements incluant contrôles brefs et facteurs
  confondants;
  \item une \textbf{évaluation} du temps d'exécution de la chaîne complète sur les 28 vaches;
  \item une \textbf{analyse complémentaire} IceTag-SLS synchronisée, explicitement interprétée comme
  exploratoire et non concluante;
  \item une \textbf{traçabilité} par tests et artefacts tabulaires reproductibles.
\end{enumerate}

La portée reste clairement délimitée: la sortie est une alerte de dégradation
comportementale à vérifier, et non une classification ou un diagnostic vétérinaire de
boiterie. Cette distinction est fondamentale et cohérente avec les recommandations de la
littérature \citep{alsaaod2019automatic}.

Le Chapitre~\ref{chapitre:systeme} définira alors le système proposé et le protocole selon lequel
il est évalué, puis le Chapitre~\ref{chapitre:donnees} décrira les données utilisées, les
transformations appliquées aux signaux capteurs, le choix des variables et le cadre
expérimental retenu. L'objectif n'est pas seulement de décrire une implémentation, mais de
montrer comment la revue qui précède se traduit en décisions méthodologiques explicites.
